\clearpage
\hypertarget{native-491-zach}{}
\label{native-491-zach}
\begin{center}
VON ZACH AN GAUSS.
\end{center}

\begin{flushright}Eisenberg, 6. April 1807.\end{flushright}

Mein hochverehrter Freund, ich bin Ihnen für die freundliche Übersendung der Broschüre von \textsc{Güssmann} unendlich verbunden, noch mehr aber für Ihre gründliche Beurteilung. Ich werde sie in das Maiheft aufnehmen, das gerade im Druck ist, und Ihren Namen dabei auf Ihren Wunsch verschweigen, so gern ich ihn genannt hätte. Aber natürlich muss und werde ich Ihren Willen und Ihre Ruhe jederzeit respektieren. Ich bin schon zufrieden, dass Sie meine Bitte angehört haben und meine Ansicht für begründet halten, dass man so böse und boshafte Menschen wie die Gefährten unseres P. \textsc{Güssmann} in ihrem wahren Wesen zeigen muss, um manchen verdienten Mann bei Ehren zu halten und ihn vor Nachteil und Schaden durch nicht hinreichend unterrichtete Personen zu bewahren. Ein so verdienstvolles und gutes Werk haben Sie, mein Verehrtester, sowohl für die gute Sache als auch für Ihre guten Freunde unternommen. Sie und alle rechtgesinnten Menschen werden Ihnen dafür ewig dankbar sein. --- --- ---
